\documentclass[10pt,a4paper]{amsart}
\usepackage[margin=19mm]{geometry}
\usepackage{amsmath,amssymb,mathtools}
\usepackage[scr=rsfs,cal=euler]{mathalfa}
\usepackage{xfrac}
\usepackage[unicode,hidelinks,bookmarks=false]{hyperref}
\newcommand{\ie}{i.e.\ }
\newcommand{\cf}{cf.\ }
\newcommand{\resp}{respectively\ }
\newcommand{\Subsection}{\subsection}
\providecommand{\nobreakdash}{\nobreak\hbox{-}}
\setlength{\emergencystretch}{2em}
\multlinegap0pt
\usepackage{amssymb}
\usepackage{mathtools}
\usepackage{caption}
\usepackage{tikz}
\usepackage[shortlabels]{enumitem}
\usepackage[all]{xy}
\newtheorem{lemma}{Lemma}
\usepackage{cleveref}
\crefname{lemma}{Lemma}{Lemmas}
\newcommand{\A}{\mathscr{A}}
\newcommand{\Dcal}{\mathcal{D}}
\newcommand{\Ecal}{\mathcal{E}}
\DeclareMathOperator{\Ext}{Ext}
\DeclareMathOperator{\Hom}{Hom}
\newcommand{\Lcal}{\mathcal{L}}
\DeclareMathOperator{\RHom}{\textbf{R}Hom}
\DeclareMathOperator{\coker}{coker}
\DeclareMathOperator{\mods}{mod}
\DeclareMathOperator{\projdim}{pdim}
\title{Exceptional versus $\tau$-exceptional sequences for the Auslander algebra of $K[x]/(x^t)$}
\author{Maximilian Kaipel}
\date{Source: arXiv:2602.13138v1 (CC BY 4.0)}
\begin{document}
\maketitle

\noindent\emph{Source and licence.} Exceptional versus $\tau$-exceptional sequences for the Auslander algebra of $K[x]/(x^t)$. Version arXiv:2602.13138v1 (CC BY 4.0), distributed by arXiv under the Creative Commons Attribution 4.0 International licence, http://creativecommons.org/licenses/by/4.0/, as recorded in the arXiv licence metadata for this exact version and shown on the abstract page https://arxiv.org/abs/2602.13138v1. Full licence notice: \texttt{licenses/arxiv-2602.13138-CC-BY-4.0.txt}.\par\smallskip

\noindent\textit{Editorial scope.} Counted unit: the article's lemma lem:excepmutdesc (source0.tex lines 562--571). The article's complete original proof of it is delivered with it (source0.tex lines 572--587, one \texttt{proof} environment of 16 lines). No mathematics is written, repaired or re-derived: the statement and the proof are the article's own lines, in the article's own notation, and no retained line is substituted or re-flowed.  Retained uncounted as the source-local context the proof needs: lines 529--529; lines 543--552.  Nothing was omitted from any retained span, so the package carries no omission record.  No retained line contains a citation, so the wrapper carries no bibliography and no bibliography entry.  No retained line contains a figure environment, an image-inclusion command or drawing code, so no image is delivered and the package declares no asset dependency.  Editorial formatting only, in addition: an AMS \texttt{amsart} preamble that restates the article's own declarations and macros used by the retained lines, namely the environments \texttt{lemma} on separate counters, together with the standard packages those lines need.  The retained environments print in this document as Lemma 1 in order of appearance, whereas the article prints the same items with the numbering of its own full document.  The complete native source of the article as distributed in the arXiv e-print for this exact version is bundled unchanged as \texttt{sources/194618/source0.tex}.
\begin{equation}\label{eq:leftmutexcep} \RHom(E,F) \otimes E \to F \to \mathcal{L}_E F \to  \quad .\end{equation}

\begin{lemma}\label{lem:completeexcepprop}
    Let $(E_t, \dots, E_1)$ be a complete exceptional sequence in $\mods \A_t$. The following hold:
    \begin{enumerate}
        \item $\dim \Hom(E_i, E_j) = 1 = \dim \Ext^1(E_i,E_j)$ for any $i>j$;
        \item $\Ext^2(E_i, E_j)=0$ for any $i>j$; 
        \item Every nonzero morphism $E_i \to E_{i-1}$ is either injective or surjective for all $i \in \{2, \dots, t\}$;
        \item If $E_2 \to E_1$ is injective, then $\projdim E_2 \leq 1$;
        \item If $f:E_i \to E_{i-1}$ is injective for $i \in \{2, \dots, t\}$, then $\dim \Ext^1(\coker f, E_i)=1$.
    \end{enumerate} 
\end{lemma}

\begin{lemma}\label{lem:excepmutdesc}
    Let $\Ecal = (E_t, \dots, E_1)$ be a complete exceptional sequence in $\mods \A_t$ and let $i \in \{2, \dots, t\}$. The following are equivalent:
    \begin{enumerate}
        \item any nonzero morphism $f: E_{i} \to E_{i-1}$ is injective;
        \item $\psi_i(\Ecal)$ is a complete exceptional sequence in $\mods \A_t$.
    \end{enumerate} 
    In this case $\Lcal_{E_i} E_{i-1} \in \mods \A_t$ is defined by the short exact sequence
    \begin{equation}\label{eq:leftmutseq} 0 \to \coker f \to \Lcal_{E_i} E_{i-1} \to E_i \to 0\end{equation}
    in $\mods \A_t$.
\end{lemma}

\begin{proof}
    This statement is unambiguous because $\dim \Hom(E_{i}, E_{i-1})=1$ by \cref{lem:completeexcepprop}(1) and therefore $\coker f_1 \simeq \coker f_2$ for any two $f_1, f_2 \in \Hom(E_i, E_{i-1})$. 

    The triangle defining $\Lcal_{E_i} E_{i-1}$, given by \cref{eq:leftmutexcep}, gives rise to a long exact sequence in cohomology
    \begin{equation}\begin{aligned}\label{eq:longexactcoho} 0 = H^{-1}(E_{i-1}) \to H^{-1}(\Lcal_{E_i} E_{i-1}) \to \Hom(E_i, E_{i-1})\otimes E_i \to E_{i-1} \to H^0(\Lcal_{E_i} E_{i-1}) \\
    \to \Ext^1(E_i, E_{i-1}) \otimes E_i \to H^1(E_{i-1}) = 0
    \end{aligned}\end{equation}
    using the fact that $E_{i-1}$ is concentrated in degree zero in $\Dcal^b(\mods \A_t)$ to get $H^{\neq 0}(E_{i-1})= 0$; the fact that $\Ext^2(E_i, E_{i-1})= 0$ by \cref{lem:completeexcepprop}(2) and the fact that $\Ext^{\geq 3}(E_i, E_{i-1}) = 0$ since $\A_t$ is of global dimension 2 to conclude that $\Lcal_{E_i} E_{i-1}$ is concentrated degrees -1 and 0. 

    Moreover, since $\dim \Hom(E_i, E_{i-1})=1 = \Ext^1(E_i, E_{i-1})$ by \cref{lem:completeexcepprop}(1), the tensor products in \cref{eq:longexactcoho} simplify to yield the exact sequence
    \begin{equation}\label{eq:simplifiedexcepdesc} 0 \to H^{-1}(\Lcal_{E_i} E_{i-1}) \to E_i \to E_{i-1} \to H^0(\Lcal_{E_i} E_{i-1}) \to E_i \to 0. \end{equation}

    Note that $E_i \to E_{i-1}$ is nonzero in \cref{eq:simplifiedexcepdesc}, so that \cref{lem:completeexcepprop}(3) implies that it is either surjective or injective. If it is injective it follows that $H^{-1}(\Lcal_{E_i} E_{i-1})=0$ and $\Lcal_{E_i} E_{i-1} \in \mods \A_t$ has the desired description, and implying that $\psi_i(\Ecal)$ is a complete exceptional sequence in $\mods \A_t$. 

    Conversely, if $\psi_i(\Ecal)$ is a complete exceptional sequence in $\mods \A_t$, then in particular $\Lcal_{E_i} E_{i-1} \in \mods \A_t$, implying $H^{-1}(\Lcal_{E_i} E_{i-1}) = 0$, consequently $f: E_i \to E_{i-1}$ is injective. This concludes the proof.
\end{proof}

\end{document}
